\documentclass[11pt]{article}

\usepackage[T1]{fontenc}
\usepackage{amsmath,amssymb,amsthm,mathtools}
\usepackage[margin=1.03in]{geometry}
\usepackage{microtype}
\usepackage{xcolor}
\usepackage{tikz}
\usetikzlibrary{arrows.meta}
\usepackage{enumitem}
\usepackage[colorlinks=true,linkcolor=blue!45!black,citecolor=blue!45!black,urlcolor=blue!45!black]{hyperref}
\hypersetup{
  pdftitle={A Noncircular Oval with Convex Unit-Tangent Iterates},
  pdfauthor={Dean Menezes}
}
\allowdisplaybreaks

\newtheorem{theorem}{Theorem}[section]
\newtheorem{proposition}[theorem]{Proposition}
\newtheorem{lemma}[theorem]{Lemma}
\newtheorem{corollary}[theorem]{Corollary}
\theoremstyle{remark}
\newtheorem{remark}[theorem]{Remark}

\newcommand{\R}{\mathbb R}
\newcommand{\Tmap}{\mathcal T}
\newcommand{\Bmap}{\mathcal B}
\newcommand{\eps}{\varepsilon}
\newcommand{\abs}[1]{\lvert #1\rvert}
\newcommand{\norm}[1]{\lVert #1\rVert}
\newcommand{\Per}{\operatorname{Per}}

\title{A Noncircular Oval with Convex Unit-Tangent Iterates}
\author{Dean Menezes}
\date{August 2026}

\begin{document}
\maketitle

\begin{abstract}
For an oriented regular plane curve \(\gamma\), let \(\Tmap\gamma\) be
obtained by moving each point one unit in the positive tangent direction.  We
prove that there is a smooth noncircular oval \(\Gamma_0\) for which every
iterate \(\Tmap^n\Gamma_0\) is an oval.  The construction begins with a
complete convex hairpin \(C\) satisfying \(\Tmap C=C+(V,0)\).  Periodizing
the sine of its steering angle gives exact centrally symmetric rear--front
pairs \(R_H\mapsto F_H\).  After their half-perimeters are matched, a direct
\(L^1\) estimate compares their intrinsic curvature measures.  A regularizing
backward-shadowing theorem converts the resulting inverse pseudo-orbit into an
exact infinite orbit.
\end{abstract}

\section{Introduction}

The unit-tangent transform moves each point of an oriented plane curve one
unit in the positive tangent direction.  Circles remain circles under this
operation.  We ask whether any noncircular oval has the same property.

Let
\[
  \gamma:\R/L\mathbb Z\longrightarrow\R^2
\]
be a smooth regular closed curve, parametrized by arclength \(s\), and let
\(\tau_\gamma=\gamma_s\) be its unit tangent.  Define
\begin{equation}\label{eq:Tdef}
  (\Tmap\gamma)(s)=\gamma(s)+\tau_\gamma(s).
\end{equation}
The definition is independent of the regular parametrization.  If the image
is regular, we retain the point correspondence and orientation, then
reparametrize by arclength before applying \(\Tmap\) again.

We orient convex curves counterclockwise, rotate the tangent counterclockwise
through \(\pi/2\) to obtain the inward normal \(\nu\), and write
\(\tau_s=k\nu\).  An \emph{oval} is a smooth embedded closed curve with
\(k>0\) everywhere.

\begin{theorem}\label{thm:main}
There is a noncircular oval \(\Gamma_0\) such that \(\Tmap^n\Gamma_0\) is an
oval for every integer \(n\geq0\).
\end{theorem}

Tabachnikov conjectured that a curve with this property must be a circle
\cite[Conjecture~4]{Tabachnikov2023}; Theorem~\ref{thm:main} answers that
question negatively.  A circle of radius \(r\) is sent to one of radius
\(\sqrt{r^2+1}\).  Our example instead becomes increasingly elongated.  Its
local model is a complete convex hairpin with parallel ends that \(\Tmap\)
carries to a translate of itself.

Suppose \(F=\Tmap R\).  We call \(R\) the \emph{rear track}, \(F\) the
\emph{front track}, and the angle between corresponding tangents the
\emph{steering angle}.  A prescribed front does not determine a rear until the
initial direction of the unit segment has been chosen.  Once that direction is
fixed, the first-order steering equation determines the rear track locally.
For a closed front, closed rear tracks are periodic solutions, or fixed points
of the bicycle monodromy; there may be none, one, or several.  In the
low-curvature class used here, there is a unique periodic solution on the
branch \(0<\delta<\pi/2\).  We denote this selected inverse by \(\Bmap\); the
notation refers only to this branch.

The proof has four steps.

\begin{enumerate}[label=\arabic*.,leftmargin=*]
\item An order-preserving operator and explicit barriers produce a complete
      strictly convex hairpin \(C\) with \(\Tmap C=C+(V,0)\).
\item The steering mass \(y=\sin\delta\) of this hairpin is a positive,
      exponentially localized pulse of mass \(\pi\).  Periodizing \(y\)
      gives exact centrally symmetric pairs \(\Tmap R_H=F_H\).
\item After matching half-perimeters, the rear curvature measure of \(R_H\)
      and the front curvature of \(F_{P(H)}\) are both compared with the same
      periodization of the isolated hairpin curvature.  Their \(L^1\)
      difference is exponentially small.
\item Curvature interpolation turns these discrepancies into summable normal
      paths.  Exact Jacobi estimates for \(\Bmap\) keep the paths in a uniform
      low-curvature class and yield a convergent sequence of terminal
      pullbacks.
\end{enumerate}

\section{One tangent step: rear and front tracks}\label{sec:bicycle}

Let $F=\Tmap R$.  Parametrize $F$ by front arclength $s$, write $x$ for rear
arclength, and choose tangent angles $\Theta,\Psi$ so that
\[
  \tau_F=(\cos\Theta,\sin\Theta),
  \qquad
  \tau_R=(\cos\Psi,\sin\Psi).
\]
Put $\delta=\Theta-\Psi$.  Resolving the front velocity in the rear Frenet
frame gives
\begin{equation}\label{eq:bicycle-s}
  \frac{dx}{ds}=\cos\delta,
  \qquad
  \frac{d\Psi}{ds}=\sin\delta.
\end{equation}
If $K=\Theta_s$ and $k=\Psi_x$ are the front and rear curvatures, then
\begin{equation}\label{eq:bicycle-curv}
  K=\delta_s+\sin\delta,
  \qquad
  k=\tan\delta.
\end{equation}
If the front is strictly convex and $\phi=\Theta$ is used as parameter, with
radius of curvature $q=1/K$, then
\begin{equation}\label{eq:bicycle-angle}
  \delta_\phi=1-q(\phi)\sin\delta.
\end{equation}
These are the standard bicycle equations; see, for example,
Levi--Tabachnikov~\cite{LeviTabachnikov}.  Equation~\eqref{eq:bicycle-angle}
also makes the preceding qualification precise: an initial value
\(\delta(\phi_0)\) determines a rear track, while a closed rear track is a
periodic solution.

\begin{lemma}[Low-curvature inverse]\label{lem:lowcurv}
Let $F$ be a smooth strictly convex closed curve with
\[
  0<K\leq\kappa<1.
\]
There is a unique closed rear track on the branch $0<\delta<\pi/2$.  It is
smooth and strictly convex, and
\begin{equation}\label{eq:rear-k-bound}
  k_R\leq\frac{\kappa}{\sqrt{1-\kappa^2}}.
\end{equation}
If $F$ is centrally symmetric, so is the selected rear track.
\end{lemma}

\begin{proof}
Set \(d_0=\arcsin\kappa\).  In \eqref{eq:bicycle-angle}, the vector field is
positive at \(0\) and nonpositive at \(d_0\).  Thus the time-\(2\pi\)
map sends \([0,d_0]\) into itself and has a fixed point.  If \(\delta_1\)
and \(\delta_2\) are two solutions in this interval, their difference
satisfies
\[
  w'=-q(\phi)\cos\xi(\phi)\,w
\]
for a value \(\xi(\phi)\) between \(\delta_1(\phi)\) and
\(\delta_2(\phi)\).  One full turn strictly decreases \(|w|\), so the
periodic solution is unique.  Its rear tangent angle is strictly increasing,
and
\[
   k_R=\tan\delta\leq\tan d_0
      =\frac{\kappa}{\sqrt{1-\kappa^2}}.
\]
Half-turn symmetry follows from uniqueness.
\end{proof}

\begin{lemma}[Convex consecutive tracks]\label{lem:strictness}
Let $\gamma$ be a smooth closed convex curve.  If $\Tmap\gamma$ is also
convex, then $\gamma$ is strictly convex.
\end{lemma}

\begin{proof}
Let $k\geq0$ be the curvature of $\gamma$ and put
\[
  u=\frac{k}{\sqrt{1+k^2}}.
\]
A Frenet calculation gives, at corresponding points,
\[
  K=u_x+u,
\]
where \(K\geq0\) is the curvature of \(\Tmap\gamma\).  Therefore
\[
   \bigl(e^x u(x)\bigr)'=e^xK(x)\geq0.
\]
If \(u(x_0)=0\) at some point, periodicity would force \(u\equiv0\),
contrary to the turning number of a closed convex curve.  Hence \(k>0\).
\end{proof}

\section{A translating hairpin}\label{sec:translator}

A \emph{hairpin} is a complete embedded strictly convex curve whose tangent
turns through \(\pi\) and whose two ends are asymptotic to parallel lines.  We
construct one satisfying \(\Tmap C=C+(V,0)\).

Parametrize the curve by tangent angle $0<\theta<\pi$ and write
\[
  \tau(\theta)=(\cos\theta,\sin\theta),
  \qquad
  C_\theta=\rho(\theta)\tau(\theta),
  \qquad
  f(\theta)=\rho(\theta)\sin\theta.
\]
If $C=(X,Z)$, then
\begin{equation}\label{eq:XZ}
  X'=f\cot\theta,
  \qquad
  Z'=f.
\end{equation}
The tangent angle of $C+\tau$ is
\[
  g(\theta)=\theta+d(\theta),
  \qquad
  d(\theta)=\arctan\frac{\sin\theta}{f(\theta)}.
\]

\begin{lemma}[Translator equation]\label{lem:translator-equation}
Suppose that $f$ is smooth on $(0,\pi)$, bounded above, and bounded below by a
constant greater than $1$.  If
\begin{equation}\label{eq:translator-integral}
  \int_\theta^{\theta+d(\theta)}f(t)\,dt=\sin\theta,
  \qquad
  d(\theta)=\arctan\frac{\sin\theta}{f(\theta)},
\end{equation}
then $g=\theta+d$ is an increasing diffeomorphism of $(0,\pi)$ and
\begin{equation}\label{eq:translator-vector}
  C(\theta)+\tau(\theta)=C(g(\theta))+(V,0)
\end{equation}
for a constant $V$.
\end{lemma}

\begin{proof}
Differentiating \eqref{eq:translator-integral} gives
\[
  f(g)g'=f+\cos\theta>0.
\]
The integral identity gives the vertical component of
\eqref{eq:translator-vector}.  For the horizontal component, put
$V(\theta)=X(\theta)+\cos\theta-X(g(\theta))$.  Using
$\tan(g-\theta)=\sin\theta/f(\theta)$ and the displayed formula for $g'$
gives $V'=0$.
\end{proof}

For a bounded measurable function
\(f:(0,\pi)\to\mathbb R\) satisfying
\[
   1<m\leq f\leq M<\infty,
\]
let
\[
   A_f(t)=\int_0^t f(u)\,du.
\]
For \(0<\theta<\pi\), define
\begin{equation}\label{eq:Df-def}
   D_f(\theta)
   =A_f^{-1}\bigl(A_f(\theta)+\sin\theta\bigr)-\theta.
\end{equation}
This is well defined because
\[
   A_f(\pi)-A_f(\theta)
      \geq m(\pi-\theta)>\sin\theta.
\]
Set
\begin{equation}\label{eq:Poperator}
   (\mathcal Pf)(\theta)=\sin\theta\cot D_f(\theta),
   \qquad 0<\theta<\pi.
\end{equation}
The endpoints \(0\) and \(\pi\) are not part of the domain.  Excluding them
allows the monotone limit to have an endpoint layer even though every finite
iterate extends continuously to \([0,\pi]\).

\begin{lemma}[Monotone translator operator]\label{lem:Pmonotone}
The map \(\mathcal P\) sends bounded profiles with
\(1<m\leq f\leq M\) to positive continuous functions on \((0,\pi)\).
It is order preserving:
\[
   f\leq h\quad\Longrightarrow\quad \mathcal Pf\leq\mathcal Ph.
\]
A bounded profile \(f\) satisfies \(f=\mathcal Pf\) on \((0,\pi)\) if
and only if it satisfies the translator equation
\eqref{eq:translator-integral} there.
\end{lemma}

\begin{proof}
The primitive \(A_f\) is continuous and strictly increasing, so
\(D_f\), and hence \(\mathcal Pf\), is continuous on the open interval.
Also
\[
   0<D_f(\theta)\leq\frac{\sin\theta}{m}<1<\frac\pi2,
\]
so \(\mathcal Pf>0\).

If \(f\leq h\), then, for fixed \(\theta\), the integral of \(h\) reaches
\(\sin\theta\) no later than the integral of \(f\).  Hence
\(D_h\leq D_f\), and the monotonicity of cotangent gives
\(\mathcal Pf\leq\mathcal Ph\).

The fixed-point equation \(f=\mathcal Pf\) is equivalent to
\[
   \tan D_f(\theta)=\frac{\sin\theta}{f(\theta)}.
\]
Since \(0<D_f<\pi/2\), this says precisely that \(D_f=d\) in
\eqref{eq:translator-integral}.  The definition of \(D_f\) now gives the
integral equation.  Conversely, that equation gives \(D_f=d\), and hence
\(f=\mathcal Pf\).
\end{proof}

Put
\[
   u_0(\theta)=\frac23(1-\cos\theta),
   \qquad
   v_-(\theta)=-\cos\theta,
   \qquad
   v_+(\theta)=\cos\theta+2,
\]
and, for \(\eps>0\), define
\begin{equation}\label{eq:barriers}
   f_\eps^\pm(\theta)
      =\eps^{-1}+u_0(\theta)+\eps v_\pm(\theta).
\end{equation}

The barriers satisfy
\[
   f_\eps^+(\theta)-f_\eps^-(\theta)
      =2\eps(1+\cos\theta)\geq0.
\]
For sufficiently small positive \(\eps\), their continuous extensions to
\([0,\pi]\) satisfy
\[
   \min_{0\leq\theta\leq\pi} f_\eps^-(\theta)
      =\eps^{-1}-\eps,
   \qquad
   \min_{0\leq\theta\leq\pi} f_\eps^+(\theta)
      =\eps^{-1}+3\eps.
\]
Thus both barriers belong to the domain of \(\mathcal P\), and
\(f_\eps^-\leq f_\eps^+\).

\begin{lemma}[Explicit barriers]\label{lem:barriers}
For all sufficiently small \(\eps>0\),
\begin{equation}\label{eq:barrier-ineq}
   f_\eps^-\leq\mathcal Pf_\eps^-
   \leq\mathcal Pf_\eps^+\leq f_\eps^+
   \qquad(0<\theta<\pi).
\end{equation}
\end{lemma}

\begin{proof}
For a profile \(f\), set
\[
   d_f=\arctan\frac{\sin\theta}{f(\theta)},
   \qquad
   \mathcal R(f)
   =\int_\theta^{\theta+d_f}f(t)\,dt-\sin\theta.
\]
The integral \(D\mapsto\int_\theta^{\theta+D}f\) is strictly increasing.
Since \(D_f\) is the value at which it equals \(\sin\theta\), while
\(f(\theta)=\sin\theta\cot d_f\), one has
\[
 \mathcal R(f)\geq0
 \iff D_f\leq d_f
 \iff \mathcal Pf(\theta)\geq f(\theta),
\]
and similarly with the inequalities reversed.  Thus the sign of
\(\mathcal R(f)\) is the sign of \(\mathcal Pf-f\).
For \(f=\eps^{-1}+u_0+\eps v\), divide by the common second-order zero
and expand in \(\eps\):
\begin{equation}\label{eq:barrier-expansion}
   \frac{2\mathcal R(f)}{d_f^2\sin\theta}
   =\eps\left(
       \frac{v'(\theta)}{\sin\theta}+\frac29\cos\theta
     \right)+O(\eps^2),
\end{equation}
where the quotient is interpreted continuously at \(0\) and \(\pi\),
and the remainder is uniform in \(\theta\).

An exact factorization controls the apparent endpoint singularities.  Each
barrier has the form
\[
   f(\theta)=A+B\cos\theta.
\]
Write \(c=\cos\theta\), \(s=\sin\theta\),
\(F=A+Bc\), and \(d=\arctan(s/F)\).  Direct integration gives
\begin{equation}\label{eq:barrier-exact-residual}
   \mathcal R(f)
   =Ad+B\{s(\cos d-1)+c\sin d\}-s.
\end{equation}
Since \(F=s\cot d\), this becomes
\begin{equation}\label{eq:barrier-factorization}
 \frac{2\mathcal R(f)}{d^2s}
 =\Phi_0(d^2)+B\Phi_1(d^2)+\frac{Bc}{F}\Phi_2(d^2),
\end{equation}
where
\[
 \Phi_0(r)=\frac{2(\sqrt r/\tan\sqrt r-1)}r,
 \qquad
 \Phi_1(r)=\frac{2(\cos\sqrt r-1)}r,
\]
\[
 \Phi_2(r)=\frac{2(\sin\sqrt r-\sqrt r)}
                  {r\tan\sqrt r}.
\]
All three functions extend analytically through \(r=0\), with
\[
 \Phi_0(0)=-\frac23,
 \qquad \Phi_1(0)=-1,
 \qquad \Phi_2(0)=-\frac13.
\]
If \(v(\theta)=\alpha+\beta c\), then
\[
 F=\eps^{-1}\{1+\eps u_0(c)+\eps^2v(c)\},
 \qquad B=-\frac23+\eps\beta.
\]
Also,
\[
 \left(\frac{s}{F}\right)^2
 =\frac{\eps^2(1-c^2)}{\{1+\eps u_0(c)+\eps^2v(c)\}^2},
 \qquad
 \frac{Bc}{F}
 =\frac{\eps Bc}{1+\eps u_0(c)+\eps^2v(c)}.
\]
Thus the right side of \eqref{eq:barrier-factorization} is smooth in
\((\eps,c)\) on a neighborhood of
\([0,\eps_0]\times[-1,1]\).  Its value at \(\eps=0\) is zero, and its
\(\eps\)-derivative there is
\[
   -\beta+\frac29c
   =\frac{v'(\theta)}{\sin\theta}+\frac29\cos\theta.
\]
Taylor's theorem on this compact rectangle proves
\eqref{eq:barrier-expansion} with a uniform remainder.

For \(v_-=-\cos\theta\), the coefficient in parentheses is
\[
   1+\frac29\cos\theta\geq\frac79.
\]
For \(v_+=\cos\theta+2\), it is
\[
   -1+\frac29\cos\theta\leq-\frac79.
\]
Thus \(\mathcal R(f_\eps^-)\geq0\) and
\(\mathcal R(f_\eps^+)\leq0\) for sufficiently small \(\eps\).
The middle inequality follows from Lemma~\ref{lem:Pmonotone}.
\end{proof}

\begin{theorem}[Translating hairpin]\label{thm:hairpin}
For every sufficiently small \(\eps>0\), equation
\eqref{eq:translator-integral} has a solution
\(f_\eps\in C^\infty(0,\pi)\) satisfying
\[
   f_\eps^-\leq f_\eps\leq f_\eps^+.
\]
The corresponding curve is a complete embedded strictly convex hairpin in a
strip of finite width, and
\[
   \Tmap C=C+(V,0)
\]
for some \(V>0\).
\end{theorem}

\begin{proof}
Start with \(f_0=f_\eps^-\) on \((0,\pi)\), and define
\(f_{n+1}=\mathcal Pf_n\).  The barrier lemma and monotonicity give
\[
   f_\eps^-\leq f_0\leq f_1\leq\cdots\leq f_\eps^+.
\]
Let \(f=\lim_n f_n\) pointwise on \((0,\pi)\), and set
\[
   A_n(t)=\int_0^t f_n(u)\,du,
   \qquad
   A(t)=\int_0^t f(u)\,du.
\]
Monotone convergence gives
\[
   \sup_{0\leq t\leq\pi}|A_n(t)-A(t)|
   \leq\int_0^\pi(f-f_n)\longrightarrow0.
\]
The primitives have a common positive lower bound for their slopes.  Their
inverses therefore converge uniformly on compact subintervals of their common
ranges.  For every \(0<\theta<\pi\),
\[
   D_{f_n}(\theta)
   \longrightarrow
   A^{-1}(A(\theta)+\sin\theta)-\theta
   =D_f(\theta).
\]
Passing to the limit in
\(f_{n+1}(\theta)=\sin\theta\cot D_{f_n}(\theta)\) gives
\(f=\mathcal Pf\) on the open interval.  Lemma~\ref{lem:Pmonotone}
therefore gives the translator equation.

The primitive \(A\) is continuous and strictly increasing; hence \(D_f\)
is continuous.  The identity \(f=\sin\theta\cot D_f\) then makes \(f\)
continuous on \((0,\pi)\).  The relation
\[
   A(\theta+D_f(\theta))-A(\theta)=\sin\theta
\]
may now be differentiated.  Successive differentiations give
\(D_f,f\in C^\infty(0,\pi)\).

No endpoint values are assigned; the construction uses only the profile on
\((0,\pi)\) and the uniform barrier bounds.

Define \(C\) by \eqref{eq:XZ}.  Its radius of curvature is
\(f/\sin\theta>0\).  Since \(Z'>0\), the curve is embedded.  The barrier
bounds imply that its arclength is infinite at both ends, that
\(X(\theta)\to-\infty\) at both ends, and that \(Z\) has finite one-sided
limits.  Hence the ends are asymptotic to parallel horizontal lines.
Lemma~\ref{lem:translator-equation} gives the translation law.  At
\(\theta=\pi/2\), one has \(g(\theta)>\theta\) and \(\cot t<0\) on the
interval between them; therefore
\[
   V=-\int_{\pi/2}^{g(\pi/2)}f(t)\cot t\,dt>0.
\]
\end{proof}

Fix such a hairpin so wide that its curvature is below \(1/20\).  Regard it
as the rear track and its translate as the front track.  Parametrize
the front by arclength $s$, let $\delta(s)$ be the steering angle, and put
\begin{equation}\label{eq:yc}
  y(s)=\sin\delta(s),
  \qquad
  c(s)=\cos\delta(s)=\sqrt{1-y(s)^2}.
\end{equation}
Choose front and rear arclength origins so that one positive function
$K_*:\R\to(0,\infty)$ is the intrinsic curvature of both translated copies.
If $x=x(s)$ is the corresponding rear arclength, then
\begin{equation}\label{eq:isolated-identities}
  x'=c,
  \qquad
  K_*(s)=y+\frac{y'}c,
  \qquad
  K_*(x(s))=\frac yc.
\end{equation}

\begin{lemma}[Hairpin pulse estimates]\label{lem:pulse}
The function \(y\) and all its derivatives decay exponentially at both
ends.  More precisely, there are \(a_-,a_+>0\) such
that, for every \(j\geq0\),
\[
  |y^{(j)}(s)|\leq C_j e^{a_-s}\quad(s\leq0),
  \qquad
  |y^{(j)}(s)|\leq C_j e^{-a_+s}\quad(s\geq0),
\]
and
\begin{equation}\label{eq:relative-y-derivatives}
   |y^{(j)}(s)|\leq D_jy(s).
\end{equation}
It also satisfies
\begin{equation}\label{eq:mass-defect}
   \int_{\R}y(s)\,ds=\pi,
   \qquad
   \Delta:=\int_{\R}(1-c(s))\,ds>0,
\end{equation}
and
\begin{equation}\label{eq:Kstar-lower}
   K_*(s)\geq b_0y(s)
\end{equation}
for some constant \(b_0>0\).
\end{lemma}

\begin{proof}
\emph{Tail bounds.}
Let \(u\) be rear arclength, let \(\theta(u)\) be the rear tangent angle,
and write
\[
   \kappa(u)=\theta_u=\frac{\sin\theta(u)}{f(\theta(u))}.
\]
The barrier enclosure gives fixed constants \(1<m<M\) with
\(m\leq f\leq M\).  The half-angle variable
\[
   r(u)=\tan\frac{\theta(u)}2
\]
satisfies the exact equation
\begin{equation}\label{eq:half-angle-growth}
   (\log r)_u
   =\frac{\theta_u}{\sin\theta}
   =\frac1{f(\theta)}.
\end{equation}
Hence \(r\), \(\sin\theta\), and \(\kappa\) have two-sided exponential
bounds at the two ends.  In particular, \(\kappa^2\) is integrable; indeed,
\[
   \int_\R\kappa(u)^2\,du
   =\int_0^\pi\frac{\sin\theta}{f(\theta)}\,d\theta<\infty.
\]

\emph{The shifted curvature equation.}
Let \(\sigma(u)\) be front arclength at the corresponding point.  Since
\[
   \sigma'(u)=\sqrt{1+\kappa(u)^2},
\]
the function \(\sigma(u)-u\) is bounded.  Choose the front and rear
arclength origins so that the intrinsic curvature of both translated copies
is the same function \(K_*\).  Then \(K_*(u)=\kappa(u)\), and the Frenet
formula for the front curvature is
\begin{equation}\label{eq:translator-curvature-identity}
   K_*(\sigma(u))
   =\frac{K_*'(u)+K_*(u)+K_*(u)^3}
          {(1+K_*(u)^2)^{3/2}}.
\end{equation}

\emph{Relative derivatives.}
We prove that
\begin{equation}\label{eq:relative-Kstar-derivatives}
   |K_*^{(j)}(u)|\leq C_jK_*(u)
   \qquad(j\geq0).
\end{equation}
Equation~\eqref{eq:half-angle-growth} first gives a quantitative
bounded-shift Harnack estimate.  For every \(D>0\), there is \(C_D\) such
that
\begin{equation}\label{eq:bounded-shift-Harnack}
 C_D^{-1}K_*(u)\leq K_*(u+h)\leq C_DK_*(u)
 \qquad(|h|\leq D).
\end{equation}
To prove this, note that \(r(u+h)/r(u)\) lies in a compact subinterval of
\((0,\infty)\), depending only on \(D,m,M\).  The function
\(2r/(1+r^2)=\sin\theta\) changes by a bounded factor under such a
multiplicative change of \(r\), and \(m\leq f\leq M\).  Since
\(\sigma(u)-u\) is bounded, \eqref{eq:bounded-shift-Harnack} gives
\[
   K_*(\sigma(u))\leq CK_*(u).
\]
Solving \eqref{eq:translator-curvature-identity} for \(K_*'\) now gives
\(|K_*'|\leq CK_*\).

We spell out the induction.  Assume
\eqref{eq:relative-Kstar-derivatives} through order \(j\).  From
\(\sigma'=\sqrt{1+K_*^2}\), every derivative \(\sigma^{(r)}\),
\(1\leq r\leq j+1\), is uniformly bounded; for \(r\geq2\) it is a finite
sum of products involving derivatives of \(K_*\) of order at most
\(r-1\).  Fa\`a di Bruno's formula gives
\[
 \partial_u^q(K_*\circ\sigma)
 =\sum_{\ell=1}^q
   K_*^{(\ell)}(\sigma(u))
   B_{q,\ell}(\sigma',\ldots,\sigma^{(q-\ell+1)}),
\]
where the \(B_{q,\ell}\) are Bell polynomials.  By the induction
hypothesis, \eqref{eq:bounded-shift-Harnack}, and the preceding bounds on
\(\sigma^{(r)}\), every term is \(O(K_*(u))\) for \(q\leq j\).

Differentiate
\[
   K_*'
   =(1+K_*^2)^{3/2}K_*\circ\sigma-K_*-K_*^3
\]
exactly \(j\) times.  Derivatives of the prefactor
\((1+K_*^2)^{3/2}\) are bounded; every term of positive derivative order
contains at least one derivative of \(K_*\) and is therefore
\(O(K_*)\).  The composed factor is \(O(K_*)\) by the previous paragraph,
and every derivative of \(K_*^3\) is also \(O(K_*)\).  Hence
\(|K_*^{(j+1)}|\leq C_{j+1}K_*\), which closes the induction.  The
exponential bounds for \(K_*\) therefore extend to every derivative.

\emph{The steering pulse.}
Let \(s=\sigma(u)\), and let \(x(s)=u\) be the inverse correspondence.
Since \(\tan\delta=K_*(u)\),
\begin{equation}\label{eq:y-from-Kstar}
   y(s)=\sin\delta(s)
       =\frac{K_*(x(s))}{\sqrt{1+K_*(x(s))^2}},
   \qquad
   c(s)=\frac1{\sqrt{1+K_*(x(s))^2}}.
\end{equation}
The derivatives of \(x=\sigma^{-1}\) are uniformly bounded and, beyond
first order, are smooth expressions in derivatives of \(K_*\).  Combining
\eqref{eq:y-from-Kstar} with
\eqref{eq:relative-Kstar-derivatives} and the chain rule gives
\eqref{eq:relative-y-derivatives} and exponential decay of every derivative
of \(y\).

\emph{Mass and the lower comparison.}
The mass identity is exact:
\[
   y(s)\,ds
   =\frac{K_*(u)}{\sqrt{1+K_*(u)^2}}
      \sqrt{1+K_*(u)^2}\,du
   =K_*(u)\,du.
\]
The rear tangent turns through \(\pi\), so \(\int_\R y=\pi\).  The
integral defining \(\Delta\) is finite because
\(1-c=O(y^2)\), and it is positive because \(y\not\equiv0\).

By \eqref{eq:relative-Kstar-derivatives}, \(|(\log K_*)'|\leq C\), and
\(|s-x(s)|\) is bounded.  The exact identity
\[
   K_*(x(s))=\frac{y(s)}{c(s)}
\]
therefore gives
\[
   K_*(s)\geq e^{-C|s-x(s)|}K_*(x(s))
       \geq b_0y(s)
\]
for a fixed \(b_0>0\), which proves \eqref{eq:Kstar-lower}.
\end{proof}

\section{Exact two-cap pairs}\label{sec:pairs}

For $H>0$, periodize the steering mass:
\begin{equation}\label{eq:YH}
  Y_H(s)=\sum_{m\in\mathbb Z}y(s-mH).
\end{equation}
For all sufficiently large \(H\), \(0<Y_H<1\).  Set
\begin{equation}\label{eq:cap-definitions}
  \delta_H=\arcsin Y_H,
  \qquad
  c_H=\sqrt{1-Y_H^2},
\end{equation}
\begin{equation}\label{eq:cap-curvatures}
  K_H=Y_H+\frac{Y_H'}{c_H},
  \qquad
  k_H(x_H(s))=\frac{Y_H(s)}{c_H(s)},
  \qquad
  x_H'=c_H.
\end{equation}

The next lemma controls periodization and tail overlap, including their
derivatives with respect to the period.

\begin{lemma}[Exponential periodization]\label{lem:periodization}
Let \(z>0\) be smooth, exponentially decaying with all derivatives, and
suppose \(|z^{(j)}|\leq C_jz\).  Put
\[
   Z_H(s)=\sum_{m\in\mathbb Z}z(s-mH),
   \qquad I_H=[-H/2,H/2].
\]
For every pair of nonnegative integers \(r,q\), there are
\(C_{r,q},\beta_{r,q}>0\) such that
\begin{equation}\label{eq:periodization-H-derivatives}
   \left\|\partial_H^q(Z_H-z)\right\|_{C^r(I_H)}
      \leq C_{r,q}e^{-\beta_{r,q}H}.
\end{equation}
Here \(s\) is held fixed when \(H\) is differentiated.  Furthermore, for
fixed \(j,k\),
\begin{equation}\label{eq:overlap-pointwise}
   \sum_{m\ne n}
   |z^{(j)}(s-mH)|\,|z^{(k)}(s-nH)|
   \leq Ce^{-\beta H}Z_H(s).
\end{equation}
This relative overlap estimate also remains valid after any fixed number of
\(H\)-derivatives.
\end{lemma}

\begin{proof}
For \(s\in I_H\) and \(m\ne0\),
\[
   |s-mH|\geq(|m|-1/2)H.
\]
Termwise differentiation gives
\[
   \partial_H^q\partial_s^r z(s-mH)
   =(-m)^q z^{(r+q)}(s-mH).
\]
The exponential tails dominate the polynomial factor \(|m|^q\), proving
\eqref{eq:periodization-H-derivatives}.

For the overlap estimate, write \(n=m+\ell\).  The two pulse centers are
\(|\ell|H\) apart.  Hence at least one of \(s-mH\) and \(s-nH\) lies at least
\(|\ell|H/2\) into a tail.  The relative derivative bounds reduce each fixed derivative to the estimate
for \(z_mz_n\).  Summing first over \(m\), then over
\(\ell\ne0\), proves the claim.  Differentiation in \(H\) introduces only
polynomial factors in the translation indices, and the same geometric series
absorbs them.
\end{proof}

\begin{lemma}[Front periodization error]\label{lem:front-periodization-error}
For a period \(L\), put
\[
   y_m(s)=y(s-mL),
   \qquad
   Y_L=\sum_m y_m,
   \qquad
   G(z)=\frac1{\sqrt{1-z^2}}.
\]
Then, for all sufficiently large \(L\),
\begin{equation}\label{eq:front-periodization-error}
 \left|K_L(s)-\sum_mK_*(s-mL)\right|
 \leq C\sum_{m\ne n}y_m(s)y_n(s)
 \leq Ce^{-\beta L}Y_L(s).
\end{equation}
\end{lemma}

\begin{proof}
The isolated identity gives
\[
   K_*(s-mL)=y_m(s)+G(y_m(s))y_m'(s),
\]
whereas
\[
   K_L=Y_L+G(Y_L)Y_L'.
\]
Therefore
\[
 K_L-\sum_mK_*(s-mL)
   =\sum_m y_m'\bigl(G(Y_L)-G(y_m)\bigr).
\]
For large \(L\), the functions \(Y_L\) and \(y_m\) take values in a
fixed interval \([0,a]\), where \(a<1\).  The function \(G\) is
uniformly Lipschitz on this interval, and \(|y_m'|\leq Dy_m\).  Hence
\[
\begin{aligned}
 \left|K_L-\sum_mK_*(s-mL)\right|
 &\leq C\sum_m y_m\sum_{n\ne m}y_n\\
 &=C\sum_{m\ne n}y_my_n.
\end{aligned}
\]
The second inequality in \eqref{eq:front-periodization-error} is
\eqref{eq:overlap-pointwise}.
\end{proof}

\begin{proposition}[Exact two-cap pairs]\label{prop:pairs}
For all sufficiently large \(H\), equations
\eqref{eq:YH}--\eqref{eq:cap-curvatures} define smooth strictly convex
centrally symmetric closed curves \(R_H,F_H\) such that
\[
   \Tmap R_H=F_H.
\]
The front perimeter is \(2H\), the rear perimeter is \(2P(H)\), where
\begin{equation}\label{eq:PH}
   P(H)=\int_0^Hc_H(s)\,ds,
\end{equation}
and
\begin{equation}\label{eq:P-asymptotic}
   P(H)=H-\Delta+O(e^{-\beta H}),
   \qquad
   P'(H)=1+O(e^{-\beta H}).
\end{equation}
If the translating hairpin is chosen sufficiently wide, one constant
\(\kappa_0<1\) bounds the curvatures of both closed curves for all
sufficiently large \(H\).
\end{proposition}

\begin{proof}
Choose \(\Theta_H\) with \(\Theta_H'=K_H\), and define the unit-speed
front by \(F_H'=\tau(\Theta_H)\).  Because \(\delta_H\) is periodic and
\(\int_0^H Y_H=\pi\),
\[
   \int_0^H K_H\,ds
   =\int_0^H(\delta_H'+Y_H)\,ds
   =\pi.
\]
The front tangent reverses after one \(H\)-period.  Hence the displacement
during the second period is the negative of the first, and \(F_H\) closes
after two periods.  Set \(\Psi_H=\Theta_H-\delta_H\) and
\[
   R_H=F_H-\tau(\Psi_H).
\]
Then \(\Psi_H'=Y_H\) and
\[
   R_H'=c_H\tau(\Psi_H),
\]
so \(R_H\) is the rear track with curvature \(k_H\), has half-perimeter
\(P(H)\), and also closes after two periods.  Both curves are centrally
symmetric.

The rear curvature is positive.  Lemma~\ref{lem:front-periodization-error}
gives
\[
   K_H(s)
   =\sum_m K_*(s-mH)+E_H(s),
   \qquad
   |E_H(s)|\leq Ce^{-\beta H}Y_H(s).
\]
Since \(K_*\geq b_0y\),
\[
   K_H(s)\geq\bigl(b_0-Ce^{-\beta H}\bigr)Y_H(s)>0
\]
for all sufficiently large \(H\).  A closed locally convex curve of turning
number one is the boundary of a strictly convex body; hence both tracks are
embedded ovals.

Set
\[
   \Phi(z)=1-\sqrt{1-z^2}.
\]
On the centered cell,
\[
   H-P(H)=\int_{-H/2}^{H/2}\Phi(Y_H(s))\,ds.
\]
Equation~\eqref{eq:periodization-H-derivatives}, smoothness of \(\Phi\) on a
fixed interval \([0,a]\) with \(a<1\), and the omitted exponential tails give
\[
   H-P(H)=\int_\R \Phi(y(s))\,ds+O(e^{-\beta H})
          =\Delta+O(e^{-\beta H}).
\]
To differentiate, hold \(s\) fixed in the centered integral:
\[
\begin{aligned}
   \frac d{dH}(H-P(H))
   ={}&\frac12\Phi(Y_H(H/2))
      +\frac12\Phi(Y_H(-H/2))\\
     &+\int_{-H/2}^{H/2}
       \Phi'(Y_H(s))\,\partial_HY_H(s)\,ds.
\end{aligned}
\]
The endpoint terms are exponentially small.  The last integral is also
\(O(e^{-\beta H})\) by \eqref{eq:periodization-H-derivatives}, after
replacing \(\beta\) by a smaller positive constant.  This yields the
derivative estimate in \eqref{eq:P-asymptotic}.

Choose the translating hairpin so wide that its intrinsic curvature is bounded
well below one.  The periodization and overlap estimates then keep both
\(K_H\) and \(k_H\) below one fixed \(\kappa_0<1\) for all sufficiently
large \(H\).
\end{proof}

\begin{lemma}[Uniform transverse width]\label{lem:width}
After a rigid motion, let the tangent angle of \(F_H\) on the centered
half-period increase from \(0\) to \(\pi\).  Its displacement perpendicular
to the asymptotic sides is
\[
   W_H=\int_{-H/2}^{H/2}\sin\Theta_H(s)\,ds.
\]
There is a constant \(C\) such that
\[
   0<W_H\leq C
\]
for all sufficiently large \(H\).
\end{lemma}

\begin{proof}
Let \(\Theta_*\) be the tangent angle of the isolated translated hairpin,
normalized by \(\Theta_*(-\infty)=0\).  The periodization estimates and the
formula for \(K_H\) give, on the centered cell,
\[
   \|K_H-K_*\|_\infty\leq Ce^{-\beta H}.
\]
Also \(\Theta_*(-H/2)=O(e^{-\beta H})\).  With
\(\Theta_H(-H/2)=0\), integration therefore gives
\[
   \|\Theta_H-\Theta_*\|_{L^\infty(I_H)}
      \leq CH e^{-\beta H}
      \leq C'e^{-\beta'H}.
\]
Consequently,
\[
\begin{aligned}
   W_H
   &\leq\int_{I_H}\sin\Theta_*(s)\,ds
       +H\|\Theta_H-\Theta_*\|_\infty\\
   &\leq\int_\R\sin\Theta_*(s)\,ds+1
\end{aligned}
\]
for all large \(H\).  The last integral is the finite strip width of the
translating hairpin.
\end{proof}

\section{Curvature-measure matching}\label{sec:matching}

The isolated rear and front are translates, so their intrinsic curvatures are
the same function \(K_*\) after independent choices of arclength origin.  Use
the front origin to fix the phase of \(y\) and its periodizations, and set
\(x_H(0)=x(0)\) for the closed rear.  For a period \(L\), write
\[
   \overline K_L(u)=\sum_{j\in\mathbb Z}K_*(u-jL).
\]
Lemma~\ref{lem:front-periodization-error} already compares the front
curvature \(K_L\) with \(\overline K_L\).  We compare the rear curvature
measure with the same profile.

\begin{lemma}[Common phase and fundamental interval]\label{lem:phase-fundamental}
Let \(I_H=[-H/2,H/2]\) and \(J_H=x_H(I_H)\).  Then \(J_H\) has length
\(P(H)\), its translates by \(P(H)\mathbb Z\) tile \(\mathbb R\), and both
the rear and front comparisons use \(\overline K_{P(H)}\) with the same
central phase.
\end{lemma}

\begin{proof}
Since \(x_H'=c_H>0\), the image \(J_H\) is an interval and
\[
 |J_H|=\int_{I_H}c_H(s)\,ds=P(H).
\]
Its translates therefore tile the line.  The choices of intrinsic origins and
\(x_H(0)=x(0)\) identify the central isolated term in both comparisons with
\(K_*\).
\end{proof}

Set \(P=P(H)\).  The infinite-time argument uses only the next estimate.

\begin{theorem}[Curvature-measure matching]\label{thm:L1match}
There are constants $C,\beta>0$ such that
\begin{equation}\label{eq:L1match}
  \int_{\R/P(H)\mathbb Z}
  \abs{k_H(u)-K_{P(H)}(u)}\,du
  \leq Ce^{-\beta H}
\end{equation}
for all sufficiently large $H$.
\end{theorem}

\begin{proof}
Put $\alpha=\min(a_-,a_+)$.  Exponential tails imply, for every
$0<\beta<\alpha$,
\begin{equation}\label{eq:convolution}
  \int_\R y(s)y(s-r)\,ds\leq C_\beta e^{-\beta\abs r}.
\end{equation}
The integrand is bounded by
\(Ce^{-\alpha(\abs s+\abs{s-r})}\).  Consequently, for pulses
$y_m(s)=y(s-mL)$,
\begin{equation}\label{eq:pairwise-L1}
  \int_{-L/2}^{L/2}\sum_{m\ne n}y_m(s)y_n(s)\,ds
  \leq Ce^{-\beta L}.
\end{equation}
This follows by writing $n=m+q$, unfolding the sum over $m$ to $\R$, and
summing \eqref{eq:convolution} over $q\ne0$.

We compare both curvatures with \(\overline K_P\).

Let $I_H=[-H/2,H/2]$, use the phase convention above, and put
$J_H=x_H(I_H)$.  The interval $J_H$ has length $P$.  Positivity and the tail
bounds give
\[
  \norm{Y_H-y}_{L^1(I_H)}+
  \norm{Y_H-y}_{L^\infty(I_H)}\leq Ce^{-\alpha H/2}.
\]
Since $c_H=\sqrt{1-Y_H^2}$ and $c=\sqrt{1-y^2}$ stay uniformly positive,
\[
  \norm{c_H-c}_{L^1(I_H)}\leq Ce^{-\alpha H/2},
  \qquad
  \norm{x_H-x}_{L^\infty(I_H)}\leq Ce^{-\alpha H/2}.
\]
The identity $K_*(x(s))=y(s)/c(s)$ and the exact measure formula
\begin{equation}\label{eq:rear-measure}
  k_H(x_H(s))\,dx_H=Y_H(s)\,ds
\end{equation}
now give
\[
\begin{aligned}
&\norm{k_H(u)\,du-\overline K_P(u)\,du}_{\mathrm{TV}}\\
&\quad=
\int_{I_H}\abs{Y_H-c_H\overline K_P(x_H)}\,ds.
\end{aligned}
\]
Split the integrand as
\[
\begin{aligned}
Y_H-c_H\overline K_P(x_H)
={}&(Y_H-y)
 +\bigl(y-c_HK_*(x_H)\bigr)\\
&-c_H\sum_{j\ne0}K_*(x_H-jP).
\end{aligned}
\]
The first term is exponentially small.  For the second, use
\(y=cK_*(x)\), boundedness of \(K_*\) and \(K_*'\), and the preceding
estimates.  For some \(\beta>0\),
\[
 \int_{I_H}|y-c_HK_*(x_H)|\,ds
 \leq CH e^{-\alpha H/2}
 \leq Ce^{-\beta H}.
\]
For the last term, change variables $u=x_H(s)$:
\[
\begin{aligned}
\int_{I_H}c_H\sum_{j\ne0}K_*(x_H-jP)\,ds
&=\int_{J_H}\sum_{j\ne0}K_*(u-jP)\,du\\
&=\int_{\R\setminus J_H}K_*(v)\,dv.
\end{aligned}
\]
Lemma~\ref{lem:phase-fundamental} gives the last equality.  We have \(x_H(s)=s+O(1)\) and
\(P=H+O(1)\).  Hence the endpoints of \(J_H\) are \(\pm H/2+O(1)\),
and the exponential tails of \(K_*\) make the last integral exponentially
small.  Therefore
\begin{equation}\label{eq:rear-common}
  \norm{k_H-\overline K_P}_{L^1(\R/P\mathbb Z)}
  \leq Ce^{-\beta H}.
\end{equation}

For the front of period \(L=P\),
Lemma~\ref{lem:front-periodization-error} and
\eqref{eq:pairwise-L1} yield
\begin{equation}\label{eq:front-common}
  \norm{K_P-\overline K_P}_{L^1(\R/P\mathbb Z)}
  \leq Ce^{-\beta P}\leq C'e^{-\beta'H}.
\end{equation}
Combining \eqref{eq:rear-common} and \eqref{eq:front-common} proves
\eqref{eq:L1match}.
\end{proof}

\section{Regularizing backward shadowing}\label{sec:shadow}

Throughout this section, curves are centrally symmetric, centered at the
origin, and supplied with a marked point and tangent direction.  These
markings fix intrinsic-arclength phase and rotation.

\begin{lemma}[Compatible markings]\label{lem:compatible-markings}
Let \(\Tmap R=F\) be a centered centrally symmetric rear--front pair, and
fix an intrinsic point of \(R\).  Given a centered marked curve \(G\), there
is a unique rotation of the whole pair for which the tangent of \(R\) at the
chosen point agrees with the marked tangent of \(G\).  This rotation preserves
\(\Tmap R=F\), all intrinsic curvature functions, and every estimate below.
A constant shift of arclength changes only the marked phase and leaves
\(W\) and the quantities \(S_j\) unchanged.
\end{lemma}

\begin{proof}
Rotate \(R\) and \(F\) by the same angle.  The unit-tangent transform
commutes with rigid motions, whereas intrinsic curvature and arclength are
rotation invariant.  A constant arclength shift is a reparametrization, so it
changes neither the geometric path nor any normal-variation norm.
\end{proof}

We use Lemma~\ref{lem:compatible-markings} whenever consecutive model curves
are compared: rotate the whole later pair so that the marked tangent of its
selected rear agrees with the marked tangent of the earlier front.

For a smooth path, write
\[
  X_t=\xi\tau+\eta\nu.
\]
The tangential component only changes the parametrization; the normal
component \(\eta\) changes the geometric curve.  Set
\[
  W(\Gamma)=\int_0^1\norm{\eta_t}_{L^1}\,dt,
  \qquad
  S_j(\Gamma)=\int_0^1\norm{\partial_s^j\eta_t}_{L^\infty}\,dt.
\]
In this section, \emph{marked geometric \(C^2\) convergence} means uniform
convergence, in one transported periodic parameter, of
\[
   X,\qquad \tau,\qquad \kappa,\qquad g=|X_u|,
\]
with the limiting speed \(g\) bounded away from zero.  Paths may be piecewise
smooth in the path parameter.  The functionals \(W,S_0,S_1,S_2\) are
additive under concatenation, and rounding finitely many path corners changes
them by an arbitrarily small amount.

Parametrize a countable concatenation on subintervals whose endpoints
accumulate only at \(t=1\).  Summability of the \(S_j\)-lengths
gives a continuous extension at \(t=1\); equivalently, the concatenation is an
absolutely continuous path in the marked geometric topology.

\begin{lemma}[Curvature interpolation]\label{lem:curv-interp}
Let \(G_0,G_1\) be marked centrally symmetric ovals with the same
half-perimeter \(L\) and the same marked tangent direction.  Suppose their
intrinsic half-periodic curvatures \(\kappa^{(0)},\kappa^{(1)}\) satisfy
\(0<\kappa^{(i)}\leq\kappa_*\) for \(i=0,1\).  There is a path
\(\Gamma\) from \(G_0\) to \(G_1\) whose curvature is at most
\(\kappa_*\) and such that
\begin{equation}\label{eq:interp-estimate}
  W(\Gamma)+S_0(\Gamma)+S_1(\Gamma)
  \leq C(1+L)^2
  \norm{\kappa^{(1)}-\kappa^{(0)}}_{L^1(0,L)}.
\end{equation}
\end{lemma}

\begin{proof}
Set \(\kappa_t=(1-t)\kappa^{(0)}+t\kappa^{(1)}\).  Let \(\theta_t(0)\) be the
common marked tangent angle, and define
\[
   \theta_t(s)=\theta_t(0)+\int_0^s\kappa_t(r)\,dr.
\]
For \(0\leq s\leq L\), put
\begin{equation}\label{eq:explicit-curvature-interpolation}
 X_t(s)=\int_0^s\tau(\theta_t(r))\,dr
       -\frac12\int_0^L\tau(\theta_t(r))\,dr,
\end{equation}
and extend by
\[
   X_t(s+L)=-X_t(s).
\]
Because \(\int_0^L\kappa_t=\pi\), this extension is smooth.  Formula
\eqref{eq:explicit-curvature-interpolation} makes closure, central symmetry,
centering, and the marked tangent phase explicit.  Its curvature is
\(\kappa_t\), so the whole path is strictly convex and has curvature at most
\(\kappa_*\).

Set \(\epsilon=\norm{\kappa^{(1)}-\kappa^{(0)}}_{L^1}\).  Then
\(\norm{\dot\theta_t}_\infty\leq\epsilon\), and differentiating
\eqref{eq:explicit-curvature-interpolation} gives
\(\norm{\dot X_t}_\infty\leq(3/2)L\epsilon\).  If
\(\dot X_t=\xi\tau+\eta\nu\), unit-speed parametrization gives
\(\dot\theta_t=\eta_s+\kappa_t\xi\).  Hence
\[
  \norm{\eta}_\infty\leq C L\epsilon,
  \qquad
  \norm{\eta_s}_\infty\leq C(1+L)\epsilon.
\]
Integrating over the full perimeter proves \eqref{eq:interp-estimate}.
\end{proof}

\begin{lemma}[Smooth dependence of the selected rear]\label{lem:path-inverse}
Let \(F_t\), \(0\leq t\leq1\), be a smooth path of strictly convex closed
fronts whose curvatures satisfy \(0<K_t\leq\widehat\kappa<1\).  Then their
selected rear tracks form a smooth path.
\end{lemma}

\begin{proof}
Parametrize each front by its tangent angle
\(\phi\in\mathbb R/2\pi\mathbb Z\), and let \(q_t=1/K_t\) be its radius of
curvature.  The selected steering angle satisfies
\[
   \delta_\phi=1-q_t(\phi)\sin\delta.
\]
Its periodic linearization is
\[
   w_\phi+q_t\cos\delta\,w=f.
\]
On the selected strip \(0\leq\delta\leq\arcsin\widehat\kappa\), the
zeroth-order coefficient is bounded below by
\[
   \widehat\kappa^{-1}\sqrt{1-\widehat\kappa^2}>0.
\]
The periodic linear operator is invertible, with a positive Green kernel and
uniform bounds.  The hyperbolic fixed point of the Poincar\'e map therefore
depends smoothly on \(t\), as does the reconstructed rear track.
\end{proof}

\begin{lemma}[Inverse Jacobi estimates]\label{lem:jacobi}
Fix $\widehat\kappa<1$.  Let $\Gamma$ be a smooth path of strictly convex
fronts of curvature at most $\widehat\kappa$, and let $\Bmap\Gamma$ be the
path of selected rears.  There are constants depending only on
$\widehat\kappa$ such that
\begin{align}
  W(\Bmap\Gamma)&\leq W(\Gamma),\label{eq:Wnonexp}\\
  S_0(\Bmap\Gamma)&\leq C_0W(\Gamma),\label{eq:S0gain}\\
  S_1(\Bmap\Gamma)&\leq C_1\bigl(W(\Gamma)+S_0(\Gamma)\bigr),\label{eq:S1gain}\\
  S_2(\Bmap\Gamma)&\leq C_2\bigl(W(\Gamma)+S_0(\Gamma)+S_1(\Gamma)\bigr).
  \label{eq:S2gain}
\end{align}
The selected inverse gains one derivative: a $C^r$ front, $r\geq2$, has a
$C^{r+1}$ selected rear.
\end{lemma}

\begin{proof}
Let $R_t$ be the rear path, parametrize $R=R_0$ by rear arclength $x$, and
write
$\dot R=\xi\tau+\eta_R\nu$.  Put $F=R+\tau$ and let $\eta_F$ be the front
normal velocity.  Since
$\dot\tau=(\eta_{R,x}+k\xi)\nu$ and
$\nu_F=-\sin\delta\,\tau+\cos\delta\,\nu$, the tangential terms cancel:
\begin{equation}\label{eq:inverse-jacobi}
  (1+\partial_x)\eta_R
  =\sec\delta\,\eta_F\circ s.
\end{equation}
On a rear circle of length $\ell$, the periodic inverse of
$1+\partial_x$ is
\[
  (\mathcal R_\ell f)(x)
  =\frac{1}{1-e^{-\ell}}
    \int_{x-\ell}^x e^{-(x-t)}f(t)\,dt.
\]
Its kernel is positive and has integral one.  Because
$dx=\cos\delta\,ds$, \eqref{eq:Wnonexp} follows.

Lemma~\ref{lem:lowcurv} bounds \(\delta\) away from \(\pi/2\) and gives a
uniform upper bound for the rear curvature.  Total turning then bounds the rear
perimeter away from zero.  Hence \(\mathcal R_\ell\) maps \(L^1\) uniformly
into \(L^\infty\), which proves \eqref{eq:S0gain}.  From
\[
  \eta_{R,x}=\sec\delta\,\eta_F\circ s-\eta_R
\]
we get \eqref{eq:S1gain}.  Differentiating once gives
\[
  \eta_{R,xx}
  =(\sec\delta)_x\eta_F\circ s
   +\sec^2\delta\,(\eta_F)_s\circ s
   -\eta_{R,x}.
\]
The identity
\(\delta_x=\sec\delta(K-\sin\delta)\) bounds the coefficient uniformly and
proves \eqref{eq:S2gain}.

For the regularity assertion, the rear tangent angle in front arclength satisfies
$\Psi_s=\sin(\Theta-\Psi)$.  If the front is $C^r$, then $\Theta$ is
$C^{r-1}$, the equation gives $\Psi\in C^r$, and one integration in rear
arclength gives a $C^{r+1}$ rear curve.
\end{proof}

\begin{lemma}[Completeness of summable normal paths]\label{lem:complete}
Let \(G_j\) be marked regular curves joined by piecewise smooth paths \(\Gamma_j\).
Suppose all intermediate curvatures are bounded by a common constant and
\[
   \sum_j\bigl(S_0(\Gamma_j)+S_1(\Gamma_j)+S_2(\Gamma_j)\bigr)<\infty.
\]
Then \(G_j\) converges in marked geometric \(C^2\) to a regular \(C^2\)
curve.
\end{lemma}

\begin{proof}
Concatenate the paths, rephasing the next path at each junction so that its
parameter agrees with the one transported from the preceding path; the
quantities \(S_j\) are unchanged by such a constant arclength shift.  If the
original velocity is \(\xi\tau+\eta\nu\) and the speed in a periodic
parameter is \(g\), solve
\[
   \phi_t(t,u)=-\frac{\xi(t,\phi(t,u))}{g(t,\phi(t,u))},
   \qquad \phi(0,u)=u.
\]
The reparametrized path has velocity \(\eta\nu\).  This transports the
chosen parameter and changes neither the geometric path nor its normal
component.  We may therefore assume
\[
   X_t=\eta\nu.
\]
If \(u\) is the transported periodic parameter and \(g=|X_u|\), then, for
our inward normal convention,
\begin{equation}\label{eq:normal-flow-identities}
   \tau_t=\eta_s\nu,
   \qquad
   \kappa_t=\eta_{ss}+\kappa^2\eta,
   \qquad
   (\log g)_t=-\kappa\eta.
\end{equation}
These formulas use the inward-normal convention fixed in the introduction.

The assumed summability makes \(X\), \(\tau\), \(\kappa\), and \(\log g\)
uniformly Cauchy along the concatenated path.  Since the total variation of
\(\log g\) is finite and every initial speed is positive, the limiting speed
\(g\) is positive.  Passing to the limit in
\[
   X_u=g\tau,
   \qquad
   \tau_u=g\kappa\nu
\]
shows that the limit is a regular \(C^2\) curve.  The transported marker and
tangent direction converge as part of the same parametrized convergence.
\end{proof}

\begin{lemma}[Stopped curvature estimate]\label{lem:stopped-curvature}
Let \(\Gamma=(C_t)_{0\leq t\leq1}\) be a smooth path written in normal
gauge, \(X_t=\eta\nu\).  Suppose
\[
 \max\kappa_{C_0}\leq\kappa_b<\kappa_e
\]
and
\[
 S_2(\Gamma)+\kappa_e^2S_0(\Gamma)
 <\kappa_e-\kappa_b.
\]
Then \(\max\kappa_{C_t}<\kappa_e\) for every \(t\).
\end{lemma}

\begin{proof}
If \(T\) were the first exit time, then on \([0,T]\) the inward-normal
flow identity would give
\[
 \|\kappa(T)-\kappa(0)\|_\infty
 \leq\int_0^T
   \bigl(\|\eta_{ss}\|_\infty+
         \kappa_e^2\|\eta\|_\infty\bigr)\,dt
 <\kappa_e-\kappa_b,
\]
contradicting \(\max\kappa(T)=\kappa_e\).
\end{proof}

\begin{lemma}[Selected inverse on the closed strip]\label{lem:closed-strip-inverse}
Let \(F\) be a regular convex \(C^2\) closed curve with
\(0\leq K\leq\widehat\kappa<1\).  There is a unique periodic steering
solution
\[
 0\leq\delta\leq\arcsin\widehat\kappa
\]
of \(\delta_s=K-\sin\delta\).  It defines a regular convex rear track.
After normalized-arclength identification of the parameter circles, this
rear depends continuously in \(C^2\) on the front.  If the front is
\(C^r\), \(r\geq2\), the rear is \(C^{r+1}\).
\end{lemma}

\begin{proof}
The vector field points into the displayed interval, so its Poincar\'e map
has a fixed point.  The difference of two solutions satisfies
\(w_s+c(s)w=0\), where
\(c(s)\geq\sqrt{1-\widehat\kappa^2}>0\); periodicity forces \(w=0\).
Also \(x_s=\cos\delta\geq\sqrt{1-\widehat\kappa^2}\), so the rear is
regular and has nonnegative curvature \(\tan\delta\).  When the perimeter
varies, write the equation in normalized arclength.  Continuous dependence for
this scalar periodic ODE then gives convergence of \(\delta\), the coordinate
change \(x\), and the reconstructed rear.  In front
arclength \(\Psi_s=\sin(\Theta-\Psi)\); the same derivative count as in
Lemma~\ref{lem:jacobi} gives the asserted regularity gain.
\end{proof}

\begin{theorem}[Regularizing backward shadowing]\label{thm:shadow}
Let \(Q_n\) be marked centrally symmetric ovals, and put
\(A_n=\Bmap Q_{n+1}\).  Assume that \(Q_n\) and \(A_n\) have the same
half-perimeter \(L_n\) and the same marked tangent direction.  Assume also
that all their curvatures are bounded by one constant \(\kappa_0<1\).  Define
\[
 e_n=(1+L_n)^2
   \|\kappa_{Q_n}-\kappa_{A_n}\|_{L^1(0,L_n)},
 \qquad
 r_n=\sum_{m=n}^{\infty}e_m.
\]
Suppose that
\begin{equation}\label{eq:shadow-hyp}
   r_0<\infty.
\end{equation}
After discarding finitely many initial indices, the terminal pullbacks
\[
   Z_n^{(N)}=\Bmap^{N-n}Q_N
\]
are defined for all \(N\geq n\) and converge in geometric \(C^2\) to
curves \(X_n\) satisfying
\[
   X_n=\Bmap X_{n+1},
   \qquad
   \Tmap X_n=X_{n+1}.
\]
Every \(X_n\) is a smooth oval.  On the retained tail, a constant
\(C=C(\kappa_0)\), independent of \(n\) and \(N\), satisfies
\begin{equation}\label{eq:shadow-model-distance}
 d_{\mathrm H}(X_n,Q_n)+|\Per(X_n)-2L_n|\leq Cr_n.
\end{equation}
Here \(d_{\mathrm H}\) denotes Hausdorff distance.  For every \(N\geq n\), there
is a concatenated normal path
\(\Gamma_{n,N}:Z_n^{(N)}\rightsquigarrow X_n\) satisfying
\begin{equation}\label{eq:shadow-terminal-tail-low}
 S_0(\Gamma_{n,N})+S_1(\Gamma_{n,N})\leq Cr_N.
\end{equation}
If \(N\geq n+1\), the same path also satisfies
\begin{equation}\label{eq:shadow-terminal-tail}
 S_2(\Gamma_{n,N})\leq Cr_N.
\end{equation}
The orbit is unique among exact inverse orbits that remain in the \(C^2\)
closures of the shrinking tubes constructed in the proof.

There is also a number \(\eta_*=\eta_*(\kappa_0)>0\) such that no initial
indices need be discarded when \(r_0\leq\eta_*\).
\end{theorem}

\begin{proof}
\emph{Defect paths and tube constants.}
Let \(C_{\mathrm{int}}=C_{\mathrm{int}}(\kappa_0)\) be the constant in
Lemma~\ref{lem:curv-interp}, put
\[
   d_n=C_{\mathrm{int}}e_n,
   \qquad
   a_n=\sum_{m=n}^\infty d_m=C_{\mathrm{int}}r_n.
\]
Lemma~\ref{lem:curv-interp} gives a defect path
\(\Lambda_n:Q_n\rightsquigarrow A_n\) satisfying
\[
   0<\kappa\leq\kappa_0,
   \qquad
   W(\Lambda_n)+S_0(\Lambda_n)+S_1(\Lambda_n)\leq d_n.
\]

Set
\[
   \bar\kappa=\frac{2\kappa_0+1}{3},
   \qquad
   \widehat\kappa=\frac{\kappa_0+2}{3},
\]
so that \(\kappa_0<\bar\kappa<\widehat\kappa<1\).  Let
\(C_0,C_1,C_2\) be the constants in Lemma~\ref{lem:jacobi} for the ceiling
\(\widehat\kappa\), and set
\begin{equation}\label{eq:tube-constants}
 M_0=\max\{1,C_0\},
 \qquad
 M_1=\max\{1,C_1(1+M_0)\}.
\end{equation}
Define
\begin{equation}\label{eq:tube-curvature-constants}
\begin{aligned}
 C_{\mathrm{tube}}
    &=C_2(1+M_0+M_1)+\widehat\kappa^2C_0,\\
 C_{\mathrm{inc}}
    &=C_2(1+M_0+M_1)+\widehat\kappa^2M_0.
\end{aligned}
\end{equation}
All these constants depend only on \(\kappa_0\).  Define \(\mathcal D_n\) to
be the set of curves that can be joined to \(Q_n\) by a piecewise smooth path
satisfying
\[
   0<\kappa\leq\widehat\kappa,
\]
\[
   W\leq a_n,
   \qquad
   S_0\leq M_0a_n,
   \qquad
   S_1\leq M_1a_n.
\]
\emph{Invariant tubes.}
Discard finitely many initial indices so that
\begin{equation}\label{eq:tube-smallness}
   C_{\mathrm{tube}}a_{n+1}<\bar\kappa-\kappa_0
\end{equation}
for every remaining \(n\).  We prove that
\begin{equation}\label{eq:tube-invariance}
   \Bmap(\mathcal D_{n+1})\subseteq\mathcal D_n.
\end{equation}

Let \(C\in\mathcal D_{n+1}\), and let \(\Gamma\) be an admissible path
from \(Q_{n+1}\) to \(C\).  By Lemma~\ref{lem:path-inverse},
\(\Bmap\Gamma\) is a smooth path.  Lemma~\ref{lem:jacobi} gives
\[
   W(\Bmap\Gamma)\leq a_{n+1},
\]
\[
   S_0(\Bmap\Gamma)\leq C_0a_{n+1},
\]
\[
   S_1(\Bmap\Gamma)
      \leq C_1(1+M_0)a_{n+1}.
\]
The second-derivative estimate gives
\[
 S_2(\Bmap\Gamma)+\widehat\kappa^2S_0(\Bmap\Gamma)
 \leq C_{\mathrm{tube}}a_{n+1}.
\]
The output starts at \(A_n\), whose curvature is at most \(\kappa_0\).
Equations \eqref{eq:tube-smallness} and
Lemma~\ref{lem:stopped-curvature} keep the entire inverse path below
\(\bar\kappa\), hence below \(\widehat\kappa\).  Concatenate it with
\(\Lambda_n\).  Since \(a_n=d_n+a_{n+1}\),
\[
   W\leq a_n,
\]
\[
   S_0\leq d_n+C_0a_{n+1}\leq M_0a_n,
\]
\[
   S_1\leq d_n+C_1(1+M_0)a_{n+1}\leq M_1a_n.
\]
These inequalities prove \eqref{eq:tube-invariance}.  They also give the
endpoint margin
\begin{equation}\label{eq:endpoint-margin}
   \max\kappa_C
      \leq\kappa_0+C_{\mathrm{tube}}a_{n+1}
      \leq\bar\kappa
\end{equation}
for every \(C\in\Bmap(\mathcal D_{n+1})\).  Every finite terminal
pullback therefore exists and lies in its corresponding tube.

\emph{Cauchy convergence.}
Fix \(n\) and compare terminal indices \(N\) and \(N+1\).  The path
\(\Lambda_N\) joins \(Q_N\) to \(A_N=\Bmap Q_{N+1}\), and
\(\Bmap^{N-n}\Lambda_N\) joins \(Z_n^{(N)}\) to
\(Z_n^{(N+1)}\).  Before any inverse is applied,
\[
   W+S_0+S_1\leq d_N.
\]
If \(N\geq n+1\), the first inverse step gives
\[
   W\leq d_N,
   \qquad
   S_0\leq M_0d_N,
   \qquad
   S_1\leq M_1d_N,
\]
\[
   S_2\leq C_2(1+M_0+M_1)d_N.
\]
The same bounds hold after every subsequent inverse step.  At each depth,
the output path starts at a finite pullback whose curvature is at most
\(\bar\kappa\).  Discard finitely many more levels so that
\[
   C_{\mathrm{inc}}d_N<\widehat\kappa-\bar\kappa
\]
for every remaining \(N\).
Lemma~\ref{lem:stopped-curvature} then keeps every propagated path below
\(\widehat\kappa\).  Consequently, for every \(N\geq n\), the increment
path satisfies
\begin{equation}\label{eq:increment-low}
   S_0+S_1\leq Cd_N,
\end{equation}
and, when \(N\geq n+1\), it also satisfies
\begin{equation}\label{eq:increment}
   S_2\leq Cd_N.
\end{equation}

Fix \(n\).  Begin the Cauchy argument with terminal indices
\(N\geq n+1\).  Equations \eqref{eq:increment-low}--\eqref{eq:increment} and
\(\sum d_N<\infty\) satisfy the hypotheses of Lemma~\ref{lem:complete};
denote the resulting limit by \(X_n\).  Concatenating the increment paths
gives \eqref{eq:shadow-terminal-tail-low}.  If \(N\geq n+1\), it also
gives \eqref{eq:shadow-terminal-tail}.  Every finite pullback at level \(n\)
is joined to \(Q_n\) by a defining tube path with
\(S_0\leq M_0a_n\) and \(W\leq a_n\).  Passing to the limit gives
\[
 d_{\mathrm H}(X_n,Q_n)\leq M_0a_n,
 \qquad
 |\Per(X_n)-2L_n|\leq\widehat\kappa a_n;
\]
the second inequality is the absolute first-variation estimate for
perimeter along a normal path.  Since \(a_n=C_{\mathrm{int}}r_n\), this proves
\eqref{eq:shadow-model-distance}.  Continuity of
\(\Tmap\gamma=\gamma+\tau_\gamma\) in \(C^1\) gives
\(\Tmap X_n=X_{n+1}\).

\emph{Convexity and regularity.}
Every finite terminal pullback is strictly convex and has turning number one.
Marked geometric \(C^2\) convergence therefore implies
\[
   \kappa_{X_n}\geq0.
\]
The unit tangents converge uniformly as maps \(S^1\to S^1\).  For all
large indices, each approximating tangent is joined to the limiting tangent by
the shorter arc of \(S^1\); hence the limiting tangent still has degree one.
Its tangent angle is nondecreasing and increases by \(2\pi\) in one circuit.

A regular closed plane curve with nonnegative curvature and turning number one
bounds a convex body: in each direction, its scalar projection has one interval
of increase and one of decrease, and therefore attains its maximum on a
supporting line.  Thus \(X_n\) is convex, and its
curvature is at most \(\bar\kappa<1\).

We establish smoothness before upgrading convexity to strict convexity.
Along every finite rear--front pair the steering angle satisfies
\(0<\delta\leq\arctan\bar\kappa\).  The curves and their unit tangents
converge, so the limiting pair has a steering angle in the closed strip
\[
   0\leq\delta\leq\arctan\bar\kappa<\frac\pi2.
\]
Lemma~\ref{lem:closed-strip-inverse} shows that \(X_n\) is the unique
regular-branch selected rear of \(X_{n+1}\).  Thus
\(X_n=\Bmap X_{n+1}\) also holds at the limit.  Iterating the regularity gain
in \(X_n=\Bmap^rX_{n+r}\) makes every \(X_n\) smooth.
Apply Lemma~\ref{lem:strictness} to the smooth consecutive convex curves
\(X_n\) and \(X_{n+1}=\Tmap X_n\).  Every \(X_n\) is strictly convex.

\emph{Uniqueness.}
Let \((Y_n)\) and \((\widetilde Y_n)\) be two exact inverse
orbits with \(Y_N,\widetilde Y_N\in\overline{\mathcal D_N}^{\,C^2}\) for
every large \(N\).  Approximate the two curves at level \(N\) by elements
of \(\mathcal D_N\), concatenate their admissible paths through \(Q_N\),
and propagate the resulting comparison path backward.  The same estimates
give \(C^2\)-distance at most \(Ca_N\) at every fixed earlier level.
First pass to the approximating limit, using
Lemma~\ref{lem:closed-strip-inverse} at each of the finitely many inverse
steps.  Then let \(N\to\infty\).
Since \(a_N\to0\), the two exact orbits coincide.

\emph{Uniform smallness.}
Only two smallness conditions require the deletion of initial indices:
\eqref{eq:tube-smallness} and the propagated-increment inequality.  Since
\(d_n\leq a_0=C_{\mathrm{int}}r_0\), the choice
\begin{equation}\label{eq:eta-star-explicit}
 \eta_*=\frac1{2C_{\mathrm{int}}}
 \min\left\{
   \frac{\bar\kappa-\kappa_0}{C_{\mathrm{tube}}},
   \frac{\widehat\kappa-\bar\kappa}{C_{\mathrm{inc}}}
 \right\}
\end{equation}
makes both inequalities valid from level zero whenever
\(r_0\leq\eta_*\).
\end{proof}

The shadowing theorem is independent of the hairpin construction: it applies
to any low-curvature inverse pseudo-orbit with summable intrinsic-curvature
defects.

\section{Proof of the main theorem}\label{sec:application}

\begin{lemma}[Large-separation threshold]\label{lem:synchronize}
For the fixed translating hairpin, there is \(H_*<\infty\) such that every
initial separation \(H_0\geq H_*\) satisfies the properties below.
\begin{enumerate}[label=\roman*.,leftmargin=*]
\item The map \(P\) is strictly increasing on the entire cap sequence,
      \(P(H)\leq H-\Delta/2\), and the recursion
      \(P(H_{n+1})=H_n\) has a unique solution with
      \[
        H_n\geq H_0+\frac{\Delta}{2}n.
      \]
\item Proposition~\ref{prop:pairs}, Lemma~\ref{lem:width}, and
      Theorem~\ref{thm:L1match} apply at every level with one common
      curvature ceiling \(\kappa_0<1\).
\item If \(e_n\) and \(r_n\) are the defect and tail quantities in
      Theorem~\ref{thm:shadow}, then
      \begin{equation}\label{eq:synchronized-tail}
       r_0\leq C(1+H_0)^2e^{-\beta' H_0}\longrightarrow0
       \qquad(H_0\to\infty).
      \end{equation}
      In particular, after increasing \(H_*\), one has
      \(r_0\leq\eta_*(\kappa_0)\), so the shadowing theorem applies from
      level zero.
\item Let \(C_{\mathrm{sh}}\) be the constant in
      \eqref{eq:shadow-model-distance} and \(C_W\) the constant in
      Lemma~\ref{lem:width}.  Increase \(H_*\), if necessary, so that
      \begin{equation}\label{eq:synchronized-noncircle}
       \frac{2H_0-C_{\mathrm{sh}}r_0}{\pi}
          > C_W+2C_{\mathrm{sh}}r_0.
      \end{equation}
\end{enumerate}
\end{lemma}

\begin{proof}
Take \(H_*\) larger than the finitely many geometric thresholds in
Proposition~\ref{prop:pairs}, Lemma~\ref{lem:width}, and
Theorem~\ref{thm:L1match}, and so large that
\[
 P'(H)\geq\frac12,
 \qquad
 |P(H)-(H-\Delta)|\leq\frac\Delta2.
\]
These inequalities make \(P\) invertible on \([H_*,\infty)\) and imply the
stated growth of \(H_n\).
After decreasing the exponent if necessary, the matching theorem gives
\[
 e_n\leq C(1+H_n)^2e^{-\beta H_{n+1}}.
\]
The bound \(H_n\geq H_0+n\Delta/2\) reduces the sum to a polynomially
weighted geometric series and proves \eqref{eq:synchronized-tail}.  Since its
right side tends to zero, the third assertion follows from the uniformity
clause in Theorem~\ref{thm:shadow}.  The left side of
\eqref{eq:synchronized-noncircle} tends to infinity, while its right side
tends to \(C_W\); one final increase of \(H_*\) proves the fourth assertion.
\end{proof}

Choose \(H_0\geq H_*\) as in Lemma~\ref{lem:synchronize}.
\emph{Construction of the pseudo-orbit.}
Define recursively
\begin{equation}\label{eq:Hrec}
  P(H_{n+1})=H_n.
\end{equation}
Then
\[
  H_{n+1}=H_n+\Delta+O(e^{-\beta H_n}),
  \qquad
  H_n\geq H_0+\frac\Delta2n.
\]
Choose compatible representatives by Lemma~\ref{lem:compatible-markings}.
Fix the centered orientation and marked phase of \(Q_0=F_{H_0}\).
Inductively, rotate the entire exact pair
\(R_{H_{n+1}}\mapsto F_{H_{n+1}}\) so that the marked tangent of
\(R_{H_{n+1}}\), at the phase used in Theorem~\ref{thm:L1match}, agrees with
the marked tangent of \(Q_n\).  Put
\[
  Q_{n+1}=F_{H_{n+1}}
\]
with this rotation, and set
\[
  A_n=\Bmap Q_{n+1}=R_{H_{n+1}}.
\]
Central symmetry keeps every representative centered at the origin, and
rotation does not alter any curvature or matching estimate.  Hence \(Q_n\)
and \(A_n\) have the same marked tangent phase and the same half-perimeter
\(H_n\).  The matching theorem gives
\[
  \norm{\kappa_{Q_n}-\kappa_{A_n}}_{L^1(0,H_n)}
  \leq Ce^{-\beta H_{n+1}}.
\]
Therefore the quantities \(e_n,r_n\) of
Theorem~\ref{thm:shadow} satisfy \eqref{eq:synchronized-tail}.  The common
curvature ceiling follows from Proposition~\ref{prop:pairs}, and the choice
of \(H_0\) ensures \(r_0\leq\eta_*(\kappa_0)\).  The shadowing theorem
therefore applies from level zero and gives a smooth exact orbit
\[
  \Tmap X_n=X_{n+1}.
\]

\emph{Excluding a circle.}
By Lemma~\ref{lem:width}, the transverse width of \(Q_0=F_{H_0}\) is
bounded uniformly for all sufficiently large \(H_0\), whereas
\(\Per(Q_0)=2H_0\).  The stated estimate \eqref{eq:shadow-model-distance} and
Lemma~\ref{lem:synchronize} show that
\[
   d_{\mathrm H}(X_0,Q_0)+|\Per(X_0)-2H_0|
      \leq Cr_0\longrightarrow0
\]
as the initial separation \(H_0\to\infty\).  A Hausdorff perturbation of
size \(d\) changes every directional width by at most \(2d\).  Therefore the
transverse width of \(X_0\) is at most \(C_W+2C_{\mathrm{sh}}r_0\), whereas a
circle with this perimeter would have width at least
\((2H_0-C_{\mathrm{sh}}r_0)/\pi\).  This contradicts \eqref{eq:synchronized-noncircle} and completes the proof of
Theorem~\ref{thm:main}.

\clearpage
\appendix
\section{Numerically generated geometry}\label{app:schematic}

The figures are numerical illustrations and play no role in the proof.  Their
coordinates were generated from the six-mode cosine profile with
\(\Lambda=25.0400224235038\).  The corresponding translation is
\(V\approx0.01945183447\); the cap-pair plot uses \(H=160\).

\input{v10_figure_coords.tex}

\begin{figure}[ht]
\centering
\begin{tikzpicture}[
  x=0.043cm,
  y=0.043cm,
  >=Latex,
  rear/.style={very thick,blue!65!black},
  front/.style={very thick,dashed,red!70!black},
  ann/.style={gray!70!black,semithick},
  tang/.style={gray!70!black,semithick,-{Latex[length=1.7mm]}}
]
  \draw[rear] plot[smooth] coordinates {\HairpinC};
  \draw[front] plot[smooth] coordinates {\HairpinTC};
  \draw[ann,<->] (-111,-40.3757)--(-111,40.3757)
    node[midway,fill=white,inner sep=1pt,font=\small] {$W$};
  \draw[tang] (56,-39.90)--(73,-39.60);
  \draw[tang] (132,0)--(132,14);
  \draw[tang] (56,39.90)--(39,39.60);
  \node[font=\small,blue!65!black] at (-72,-31) {$C$};
  \node[font=\small,red!70!black] at (-58,31) {$\Tmap C=C+(V,0)$};
  \draw[ann,-{Latex[length=1.7mm]}] (-12,-52)--(34,-52);
  \node[font=\small,gray!70!black] at (11,-59) {translation by $(V,0)$};
\end{tikzpicture}
\caption{The sampled hairpin profile and its horizontal translate by
\(V\approx0.01945183447\).  The solid curve is \(C\); the dashed curve is
\(C+(V,0)\).  The arrows show the positive tangent direction, and \(W\)
marks the strip width.}
\label{fig:hairpin-numerical}
\end{figure}

\begin{figure}[ht]
\centering
\begin{tikzpicture}[
  x=0.084cm,
  y=0.084cm,
  >=Latex,
  rear/.style={very thick,blue!65!black},
  front/.style={very thick,dashed,red!70!black},
  ann/.style={gray!70!black,semithick},
  tang/.style={gray!70!black,semithick,-{Latex[length=1.7mm]}}
]
  \draw[rear] plot[smooth cycle] coordinates {\CapRear};
  \draw[front] plot[smooth cycle] coordinates {\CapFront};
  \draw[ann,<->] (-33,43)--(33,43)
    node[midway,fill=white,inner sep=1pt,font=\small]
    {large separation parameter $H$};
  \draw[tang] (-62,0)--(-62,8);
  \draw[tang] (0,36)--(10,36);
  \draw[tang] (62,0)--(62,-8);
  \draw[tang] (0,-36)--(-10,-36);
  \node[font=\small] at (-56,5) {cap};
  \node[font=\small] at (56,5) {cap};
  \node[font=\small,blue!65!black] at (-28,-27) {$R_H$};
  \node[font=\small,red!70!black] at (19,27) {$F_H=\Tmap R_H$};
\end{tikzpicture}
\caption{A numerical reconstruction of the periodized rear--front pair at
\(H=160\).  The solid curve is \(R_H\); the dashed curve is
\(F_H=\Tmap R_H\), up to numerical quadrature error.  The two labeled ends
are the cap regions.}
\label{fig:cappair-numerical}
\end{figure}

\clearpage
\begin{thebibliography}{9}

\bibitem{LeviTabachnikov}
M.~Levi and S.~Tabachnikov,
\newblock On bicycle tire tracks geometry, hatchet planimeter, Menzin's
conjecture, and oscillation of unicycle tracks,
\newblock \emph{Experimental Mathematics} \textbf{18} (2009), no.~2,
173--186,
\newblock \href{https://doi.org/10.1080/10586458.2009.10128894}
{doi:10.1080/10586458.2009.10128894}.

\bibitem{Tabachnikov2023}
S.~Tabachnikov,
\newblock Iterating skew evolutes and skew involutes: a linear analog of
bicycle kinematics,
\newblock \emph{Moscow Mathematical Journal} \textbf{23} (2023), no.~4,
625--640,
\newblock \href{https://doi.org/10.17323/1609-4514-2023-23-4-625-640}
{doi:10.17323/1609-4514-2023-23-4-625-640}.

\end{thebibliography}

\end{document}
